% ============================================================
% Section 6: Conclusion
% ============================================================

\section{Conclusion}
\label{sec:conclusion}

We presented DocVerify, a multi-task CNN for joint detection and localization of
forged identity documents, paired with a multi-objective Pareto HPO framework
that selects hyperparameter configurations by optimizing PR-AUC and Dice
simultaneously without fixing their relative importance in advance.

Our main findings are: (i) joint multitask training with Pareto-optimized loss
weights achieves significantly higher Dice than single-task variants and
equal-weight joint training (blind test, 30 seeds, $p = 1.3 \times 10^{-4}$,
$d = 1.01$) while maintaining competitive PR-AUC; (ii) under a rigorous 10-fold
nested cross-validation protocol, DocVerify achieves PR-AUC $= 0.9921 \pm
0.0058$ and Dice $= 0.856 \pm 0.030$, demonstrating consistent generalization
across data partitions; and (iii) training exclusively on FantasyID without any external pre-training,
DocVerify achieves results broadly competitive with top DeepID 2025 Challenge
participants (F1@0.5 $= 0.969 \pm 0.014$; F1$_{\text{loc}} = 0.807 \pm 0.096$
on our internal holdout), though direct comparison is indicative given that
challenge systems were evaluated on a different official test set.

Future work includes cross-domain evaluation on additional document datasets,
integration of forgery-artifact-centric features (noise fingerprints, frequency
anomalies~\citep{cozzolino2020noiseprint}) for domain-invariant representations, and
foundation model backbones~\citep{fakeidet2025} for richer pre-training.

\section*{Acknowledgments}
This work was supported by the Instituto Nacional de Ciberseguridad (INCIBE,
Spain's National Cybersecurity Institute) and the INCIBE Chair in Cybersecurity
and Artificial Intelligence at the University of Oviedo, Spain. The author thanks
the Idiap Research Institute for releasing the FantasyID dataset.
